
%(BEGIN_QUESTION)
% Copyright 2006, Tony R. Kuphaldt, released under the Creative Commons Attribution License (v 1.0)
% This means you may do almost anything with this work of mine, so long as you give me proper credit

Anta at vi måler strømningsraten til en gass med en turbinflowmåler. Dette er en enkel turbinflowmåler, med én turbin som roterer fritt og genererer elektroniske pulser via en «pick-up»-spole som registrerer når turbinbladene passerer.

\vskip 10pt

Hvis tettheten til denne gassen plutselig øker uten at volumstrømmen endres, vil turbinhastigheten øke, avta eller forbli den samme?

\underbar{file i00730}
%(END_QUESTION)





%(BEGIN_ANSWER)


%(END_ANSWER)





%(BEGIN_NOTES)

Turbinen vil fortsette å rotere med samme hastighet, fordi dette er en direkte funksjon av fluidets {\it hastighet} forbi turbinen, og ikke av fluidets tetthet.

\vskip 10pt

Den eneste gangen fluidtetthet vil ha betydning for en turbin, er dersom turbinen ikke kan rotere fritt. Fluidtetthet påvirker direkte fluidmomentum, som igjen påvirker {\it impulsen} fra fluidet som treffer et turbinblad som gjør motstand. Ved normal drift av turbinmåler roterer turbinen imidlertid svært fritt, og «motstår» derfor ikke fluidet som passerer. Dermed blir tetthet ikke en relevant faktor.

Et viktig unntak fra denne regelen er designet med {\it twin-turbine} massestrømsmåler, som bevisst lar to turbiner arbeide mot hverandre ved å gi dem ulike bladvinkler og koble dem via en fjær til samme aksel. Her brukes differensielt dreiemoment som indikator på fluidmomentum, som er en direkte funksjon av {\it masse}strømningsrate.

%INDEX% Måling, flow: turbin

%(END_NOTES)


